\begin{textsample}{POS Dim 3 – go – Score 26.00 – go/go\_inf\_19.txt}  \label{ex:f3_pos_044}
3.3.2.
Organização da Ação Pedagógica Ao professor de Educação \textbf{Infantil} está posto o desafio de se permitir aprender e desenvolver na convivência \textbf{diária} com as \textbf{crianças} e ter clareza de que seus próprios repertórios irão influenciar de maneira direta na ampliação dos repertórios das \textbf{crianças}.
Nesse sentido, é necessário que rompa com seus próprios estereótipos, desenvolva e/ou amplie sua capacidade expressiva, \textbf{brinque} com diferentes materiais, suportes, instrumentos musicais, ouça muitas e diferentes músicas, contemple obras de arte, visite diversificados espaços de cultura, assista espetáculos, leia, se permita experimentar.
Também é preciso fazer mediações e intervenções que oportunizem a experimentação, a expressão \textbf{infantil}, bem como a apropriação por parte da \textbf{criança}, do patrimônio artístico e cultural produzido pela humanidade.
Ao professor também cabe ampliar os repertórios das \textbf{crianças}, o que pressupõe saber quais repertórios elas têm, conhecer sua realidade, o que lhes \textbf{interessa}, o que lhes chama a atenção.
É preciso então saber : Que músicas as \textbf{crianças} conhecem?
Quais ouvem com maior frequência?
Quais são seus \textbf{hábitos} e costumes?
Elas já visitaram algum espaço de cultura?
Qual?
Há algum artista em suas famílias?
Participam de alguma manifestação cultural - congada, folia de reis, capoeira, banda de música, coral, grupo de dança, de teatro ou de circo?
A comunidade da qual a instituição faz parte pode e deve ser integrada ao processo de ampliação de repertório, convidando as pessoas para socializarem seus talentos na instituição, com as \textbf{crianças} e os \textbf{adultos}.
Elas podem \textbf{contar} causos e histórias, cantar, tocar instrumentos musicais, ensinar um artesanato típico daquele ou de outro lugar, além de tantas outras coisas.
Se nessa comunidade houver algum grupo de teatro ou de dança, que ele seja convidado a se apresentar para as \textbf{crianças} na instituição, oportunizando a elas a observação do processo de montagem dos diversos elementos que compõem o espetáculo : cenário, iluminação, figurino, trilha sonora, o ensaio dos movimentos e das falas.
Ou seja, que elas possam perceber todas as ações que culminam no espetáculo que é mostrado ao público.
Conduzidos assim, o teatro e a dança, nas instituições de Educação \textbf{Infantil}, não serão resumidos a práticas pedagógicas baseadas na reprodução de uma gestualidade sem sentido para as \textbf{crianças}, tais como os ensaios enfadonhos ; dançar tentando reproduzir a letras das músicas na coreografia ; no ensaio de passos ou ainda na memorização de textos que tenham uma apresentação como único objetivo.
Às \textbf{crianças} será dada a oportunidade de viver a arte na perspectiva de processo, criativo e autoral, e não enquanto busca de um produto a ser entregue, via de regra, aos \textbf{adultos} e que tão comumente deixam as \textbf{crianças} constrangidas, quando não, assustadas.
Outro aspecto a ser pensado e cuidado pelos profissionais das instituições diz respeito ao repertório musical que é usado no trabalho com as \textbf{crianças}.
Há que se fugir da prática de oferecer a elas músicas e cantigas que já conhecem, usando o argumento de que elas gostam.
É preciso oferecer o novo, o inusitado, o surpreendente.
Também é preciso lançar mão de recursos tecnológicos para fazer chegar às \textbf{crianças} outras manifestações artísticas e culturais às quais não seja possível garantir o acesso direto, em especial àquelas de reconhecido valor universal, regional ou local.
Assim, pode -se : - exibir filmes ou slides de obras arquitetônicas, \textbf{pinturas} e esculturas famosas ou não ; - expor gravuras e desenhos, além de fotografias de obras de arte ou de lugares significativos ; - apresentar vídeos de trechos de espetáculos de dança, de circo, de teatro ; - executar músicas de diferentes estilos e épocas usando televisão, datashow, aparelhos de CD etc.
Além disso, o trabalho desenvolvido pelas instituições deve se preocupar com a organização dos espaços, sua diversificação e também as formas como são usados.
Tanto se pode garantir que as \textbf{crianças} usem os espaços disponíveis na própria instituição : a sala, os corredores, os \textbf{pátios}, os vãos, as áreas cobertas, quanto possam extrapolar as fronteiras desta.
Que \textbf{brinquem} em \textbf{pátios}, quintais, bosques, jardins, praças, tenham acesso a outros espaços significativos, pois a convivência das \textbf{crianças} com as diferentes manifestações artísticas e culturais não pode ficar restrita ao ambiente e à convivência com as pessoas da instituição.
É preciso ir além do espaço circundante, garantir -lhes o acesso a outras pessoas, outras etnias, de diferentes grupos etários, de diferentes grupos sociais e a espaços culturais diversificados : monumentos, equipamentos públicos, museus, circos, teatros, cinemas, laboratórios, planetários, praças, jardins, parques, bibliotecas e outros lugares importantes da cidade, diversos shows e espetáculos, entre tantas outras possibilidades.
Também é preciso garantir que os \textbf{bebês} participem dos passeios e outras ações realizadas fora da instituição.
A instituição é uma oportunidade singular de contato com a arte, em especial na Educação \textbf{Infantil}.
Ela deve possibilitar que as \textbf{crianças} pisem o \textbf{chão} de mármore dos museus, que possam estar no mesmo ambiente que esculturas, quadros, instalações.
Que possam apreciar espetáculos de circo, de música, de dança e de teatro, experimentar as cadeiras estofadas, visitar bibliotecas e exposições, ouvir e cantar músicas, apreciar o silêncio, os \textbf{sons} da natureza e também de seu entorno, \textbf{manusear} livros e instrumentos musicais.
A instituição de Educação \textbf{Infantil} precisa ter, em seus diferentes espaços, as marcas de autoria das \textbf{crianças} que ali habitam.
É importante que as produções delas estejam intencionalmente expostas nos mais variados lugares, dentro e fora das salas, \textbf{contando} as histórias vividas ali por \textbf{adultos} e \textbf{crianças}.
Ser espaços de vida, marcados de histórias, \textbf{cheios} de memórias, \textbf{afetos}, aprendizagens e conquistas.
Podem ser expostos desenhos, modelagens, colagens, \textbf{pinturas}, instalações, fotografias, textos, \textbf{brinquedos}, objetos que expressem o que está sendo vivido por aquelas pessoas, naquele lugar, naquele tempo.
Periodicamente, fazer a substituição do que está exposto, não permitindo a permanência de materiais desgastados e ou estragados, assim como evitando que os observadores percam o interesse pelo espaço e o condicionamento do olhar sobre o cotidiano.
A organização dos tempos deve ser objeto de \textbf{cuidado} constante dos professores quando do planejamento e desenvolvimento de seu trabalho na instituição junto às \textbf{crianças}.
O trabalho deve ser constante, cotidiano, oferecendo às \textbf{crianças} o contato regular com a música, o teatro, o desenho, a fotografia, a \textbf{pintura}, a modelagem, a colagem, a dança, a expressão de uma maneira geral.
Contudo, não se pode abrir mão do novo, do surpreendente, do inusitado, que podem encantar e mobilizar as \textbf{crianças} e os \textbf{adultos} envolvidos no processo.
Se a \textbf{pintura} acontece uma ou duas vezes ao ano, o contato com a \textbf{tinta} pode causar tamanha euforia nas \textbf{crianças} que a atividade pode se perder nesse contato.
A tônica precisa ser a constância de oportunidade para vivências, experimentação com diferentes possibilidades, oportunizando descoberta, surpresa, familiaridade com os materiais e expressão.
Na organização das \textbf{crianças} há que se oportunizar momentos em que elas trabalhem em duplas, trios, \textbf{pequenos} e grandes grupos, além de momentos em que elas atuem sozinhas e, em outros, coletivamente.
É preciso considerar o que e onde será realizado, que materiais serão disponibilizados e o tempo disponível.
As \textbf{crianças} também podem e devem ser consultadas no que diz respeito a forma como vão se agrupar, ou não, para a realização de determinada proposta.
Em cada tipo de organização das \textbf{crianças} são oportunizadas a elas interações específicas, possibilitando ainda o exercício da autoria, tanto individual quanto coletiva, compartilhada.
No que diz respeito a seleção e oferta de material, a premissa básica é a variedade, a diversificação.
Oferecer objetos sonoros, botões, sementes, terra, serragem, tesoura, potes.
Escolher os suportes para as \textbf{tintas}, os pincéis, os rolos.
Disponibilizar \textbf{tintas} – artesanais ou industrializadas – canetas hidrocor, \textbf{lápis} preto e de cor, instrumentos musicais, papeis, telas, carvão, giz pastel, giz para quadro, rolos, fitas adesivas, linhas, barbantes, cola, água, \textbf{areia}.
Oportunizar diferentes suportes para os rabiscos, as marcas, os desenhos : \textbf{paredes}, livros, mesas, papelão, cadeiras, o próprio corpo e os suportes digitais.
No caso dos \textbf{bebês}, pensar em materiais que resistam a sua investigação, que não ofereçam risco a eles e tornem possível desenvolver o que foi proposto.
Por exemplo, fixar o papel que será rabiscado na mesa ou no \textbf{chão}, com fita adesiva, para que não deslize enquanto a \textbf{criança} faz sua produção.
Às \textbf{crianças} deve ser garantida a oportunidade de experimentar, interagir, criar e aprender na relação com diferentes tipos de materiais.
Que elas possam roçar os materiais em seu próprio corpo, provar, furar, amassar, rasgar, molhar, colar, dobrar, em movimento oportunizado, permitido, incentivado.
Muitas vezes as \textbf{crianças} querem tão somente realizar experimentações em superfícies, objetos, materiais ( CUNHA, 2012 ). \textbf{Brincar}, arrastar, empurrar, sacudir, transformar e descobrir o que é possível fazer com pincel seco ou \textbf{cheio} de \textbf{tinta}, uma flauta, um chocalho, uma lata \textbf{cheia} de sementes, uma câmera ; com bolinhas-cravo, plástico-bolha, roupas, acessórios, fantasias, tecidos grandes e \textbf{pequenos}, caixas de papelão ou de \textbf{plástico}.
Para fazer mediações que efetivamente contribuam para as aprendizagens e o desenvolvimento das \textbf{crianças}, o professor precisa conhecê -las profundamente.
Observar as \textbf{brincadeiras} e as conversas ajuda a conhecer os costumes, as preferências e modos como as \textbf{crianças} pensam.
Cabe ao professor, a partir de então, criar oportunidades para que as \textbf{crianças} experimentem cores, \textbf{sons}, traços e formas, fugindo de estereótipos, vendo com sensibilidade e pensando sobre o mundo.
Garantir a seleção e oferta de materiais diversificados, desafiando as \textbf{crianças} a fazer usos inusitados deles, permitir suas experimentações, promover a contemplação e a produção de imagens, movimentos, \textbf{sons} e músicas.
Mostrar diferentes representações de um mesmo objeto, na instituição ou fora dela, analisar os traçados usados pelos artistas em seus desenhos, para que as \textbf{crianças} ampliem seu repertório visual e percebam que não existe uma única forma de desenhar, por exemplo.
Fazer perguntas que ativem a memória e agucem a criatividade : - quando você ouve esta música se lembra de alguma coisa?
De quê? - o que você pode fazer com esta fita adesiva, já que não há outra coisa para colar com ela? - é possível pintar sem usar \textbf{tintas} compradas prontas? - que tipo de \textbf{som} se pode produzir usando esses objetos? - você consegue dançar diferente de seu \textbf{colega}? - como se pode representar o personagem dessa história? - o que pode ser usado para construir o cenário?
Além de oferecer materiais e deixar fazer, é preciso observar e dialogar com a \textbf{criança}, observar seus processos e produtos.
Incentivar sua autonomia e investigação, promovendo seu protagonismo criador.
Acreditar na \textbf{criança} e na sua capacidade é fundamental.
As mediações dos professores devem ser na perspectiva de ampliar o modo de ver, experimentar, registrar e imaginar o mundo, desenvolvendo possibilidades de criação, crítica, estesia, expressão, fruição e reflexão sobre traços, \textbf{sons}, cores e formas.

% matched lemmas: adulto, afeto, areia, bebê, brincadeira, brincar, brinquedo, cheio, chão, colega, contar, criança, cuidado, diário, hábito, infantil, interessar, lápis, manusear, parede, pequeno, pintura, plástico, pátio, som, tinta
\end{textsample}
